\documentclass[10pt,letterpaper]{article}
\usepackage{cornell-notes}

\title{Chapter 27: The Backtrack Method}
\author{}
\date{}

\setDocTitle{Chapter 27: The Backtrack Method}
\setDocSubtitle{Combinatorial Algorithms}
\setDocVersion{v1.0}
\setDocStatus{Study Companion}
\setDocOwner{Combinatorial Algorithms Cornell Notes}
\setDocAuthor{}
\setDocDate{\today}
\setDocSummary{Cornell-format study and review notes for Chapter 27: The Backtrack Method.}

\setCornellCollection{Combinatorial Algorithms Cornell Notes}
\setCornellUnitType{Chapter}
\setCornellUnitNumber{27}
\setCornellUnitTitle{The Backtrack Method}
\setCornellCompanionLabel{Study and review companion}

\hypersetup{pdftitle={Chapter 27: The Backtrack Method},pdfsubject={Cornell notes},pdfauthor={}}

\begin{document}
\maketitle
\begin{tcolorbox}[colback=Paper,colframe=Ink]{\small\bfseries\color{Teal}CORNELL NOTES}\par{\LARGE\bfseries Chapter 27: The Backtrack Method (BACKTR)}\par\medskip\textbf{Topic:} A reusable supervisor for exhaustive constrained search\hfill\textbf{Date:} \today\par\textbf{Study purpose:} Separate depth-first control from problem-specific candidate generation.\end{tcolorbox}
\begin{tcolorbox}[colback=Teal!5,colframe=Teal,title=Essential question]How does a stack-based supervisor enumerate every feasible construction while allowing different candidate rules to supply the problem knowledge?\end{tcolorbox}
\begin{tcolorbox}[colback=white,colframe=Mist,title=Learning objectives]\begin{itemize}\item Explain the MAIN--CANDTE--BACKTR protocol.\item Trace the packed candidate stack.\item State the partial-solution invariant.\item Specialize the method to graph problems.\item Identify pruning, restoration, and randomness pitfalls.\end{itemize}\textbf{Prerequisites:} depth-first search, graphs, stacks, partial solutions.\end{tcolorbox}
\section*{Cornell notes}
\begin{longtable}{@{}Q!{\color{Mist}\vrule width 1pt}A@{}}\cornellhead
\cue{What is being separated?}{The supervisor handles depth-first traversal, choice, and retreat. A problem-specific routine examines the current partial vector and writes all legal candidates for its next position. This interface lets the same control logic solve many finite construction problems.}
\cue{What do the status values mean?}{The index/status variable reports: $0$, initialize a search; $1$, a complete vector is available to the caller; $2$, candidates are needed for the next position; $3$, the search is exhausted. The caller alternates between interpreting solutions and supplying candidate lists.}
\cue{Packed STACK layout}{For each vector position, STACK stores that position's candidates followed by their current count. $K$ is the current position and $M$ is the last occupied stack location. Choosing a candidate decrements its list count; exhaustion removes that block and retreats to the preceding position.}
\cue{Core invariant}{Before position $K+1$ is extended, $X(1),\ldots,X(K)$ is a feasible partial solution. STACK retains exactly the untried legal alternatives needed to resume each active depth. State changed by a choice must be restorable on retreat.}
\cue{When is a branch complete?}{After accepting a candidate, if $K=L$ the vector has the required length and status $1$ is returned. Otherwise the routine requests candidates for $K+1$. A zero-length candidate list immediately forces backtracking.}
\cue{Graph coloring example}{COLVRT proposes colors absent from already colored neighbors. Fixing the first vertex to color $1$ removes equivalent solutions caused only by permuting color names. Early conflict tests prune long before a full coloring is formed.}
\cue{Circuit examples}{For an Euler circuit, candidates are unused edges incident with the current terminal vertex; feasibility depends on edge use. For a Hamilton circuit, candidates are unused adjacent vertices, with closing-edge and reversal/rotation symmetry conditions checked at the appropriate depth.}
\cue{Spanning-tree example}{SPNTRE builds a set of edges that remains acyclic and can still connect all vertices. Candidate generation must reject cycle-forming edges and may use a spanning-forest representation to test whether two endpoints are already connected.}
\cue{Is random choice uniform?}{No. Uniformly selecting one legal candidate at each depth weights a completed object by the reciprocal of the branching factors along its path. Unless all relevant paths have equal products, the resulting distribution over completed objects is biased.}
\cue{Engineering checks}{Bound the stack for worst-case candidate blocks; generate candidates deterministically when reproducibility matters; undo every incremental mark; distinguish ``no candidate'' from ``solution complete''; and encode symmetry rules without excluding genuinely distinct objects.}
\cue{Retrieval practice}{\begin{enumerate}\item What does status $2$ request?\item Where is a candidate count stored?\item What must remain true of the partial vector?\item Why fix the first color?\item Why is local random choice usually biased?\end{enumerate}}
\end{longtable}
\begin{tcolorbox}[colback=Ink!4,colframe=Ink,title=Cornell summary]
BACKTR is a reusable depth-first supervisor. It maintains a partial vector and packed stacks of untried candidates, requests problem-specific candidates, returns completed vectors, and retreats when a list is empty. Correctness rests on a feasible-partial-solution invariant and complete restoration of mutable state. Coloring, Euler circuits, Hamilton circuits, and spanning trees differ mainly in their candidate predicates and pruning rules. The method enumerates systematically, but choosing a locally uniform random branch does not generally sample completed objects uniformly.
\end{tcolorbox}
\end{document}
